\documentclass[10pt,a4paper]{amsart}
\usepackage[margin=19mm]{geometry}
\usepackage{amsmath,amssymb,mathtools}
\usepackage[scr=rsfs,cal=euler]{mathalfa}
\usepackage{xfrac}
\usepackage[unicode,hidelinks,bookmarks=false]{hyperref}
\newcommand{\ie}{i.e.\ }
\newcommand{\cf}{cf.\ }
\newcommand{\resp}{respectively\ }
\newcommand{\Subsection}{\subsection}
\providecommand{\nobreakdash}{\nobreak\hbox{-}}
\setlength{\emergencystretch}{2em}
\multlinegap0pt
\usepackage[shortlabels]{enumitem}
\usepackage{amssymb}
\usepackage{mathtools}
\newtheorem{theorem}{Theorem}
\newcommand{\dt}{\partial_t}
\newcommand{\dtau}{\partial_\tau}
\newcommand{\dv}{\mathrm{div}\, }
\newcommand{\idx}{\ d\vec x}
\newcommand{\mrm}[1]{\mathrm{#1}}
\newcommand{\norm}[2]{\Vert #1 \Vert_{#2}}
\renewcommand{\vec}[1]{{\bf #1}}
\title{Large friction limit of the almost pressureless Euler-Poisson system}
\author{Xin Liu}
\date{Source: arXiv:2603.18346v1 (CC BY 4.0)}
\begin{document}
\maketitle

\noindent\emph{Source and licence.} Large friction limit of the almost pressureless Euler-Poisson system. Version arXiv:2603.18346v1 (CC BY 4.0), distributed by arXiv under the Creative Commons Attribution 4.0 International licence, http://creativecommons.org/licenses/by/4.0/, as recorded in the arXiv licence metadata for this exact version and shown on the abstract page https://arxiv.org/abs/2603.18346v1. Full licence notice: \texttt{licenses/arxiv-2603.18346-CC-BY-4.0.txt}.\par\smallskip

\noindent\textit{Editorial scope.} Counted unit: the article's theorem thm:limit (source0.tex lines 315--330). The article's complete original proof of it is delivered with it (source0.tex lines 332--334, one \texttt{proof} environment of 3 lines). No mathematics is written, repaired or re-derived: the statement and the proof are the article's own lines, in the article's own notation, and no retained line is substituted or re-flowed.  Retained uncounted as the source-local context the proof needs: lines 76--94; lines 98--114; lines 118--136; lines 211--235; lines 245--248; lines 257--260; lines 275--303; lines 419--422; lines 1040--1056.  Nothing was omitted from any retained span, so the package carries no omission record.  No retained line contains a citation, so the wrapper carries no bibliography and no bibliography entry.  No retained line contains a figure environment, an image-inclusion command or drawing code, so no image is delivered and the package declares no asset dependency.  Editorial formatting only, in addition: an AMS \texttt{amsart} preamble that restates the article's own declarations and macros used by the retained lines, namely the environments \texttt{theorem} on separate counters, together with the standard packages those lines need.  The retained environments print in this document as Theorem 1 in order of appearance, whereas the article prints the same items with the numbering of its own full document.  The complete native source of the article as distributed in the arXiv e-print for this exact version is bundled unchanged as \texttt{sources/196688/source0.tex}.
Let $ \rho \in \mathbb R $, $\vec u \in \mathbb R^3$, and $ \phi \in \mathbb R $ be the density, velocity, and the electric potential, respectively. 
The almost pressureless Euler-Poisson (EP) system with large friction is given by the following system of equations: 
\begin{subequations}
        % \makeatletter
        % \def\@currentlabel{EP}
        % \makeatother
        % \renewcommand{\theequation}{EP.\arabic{equation}}
        \label{sys:PLEP}
    \begin{align}
        \label{eq:PLEP-1}
        \dt \rho + \dv (\rho \vec u) & =0, & \text{in} \ \Omega, \\
        \label{eq:PLEP-2} 
        \dt (\rho \vec u) + \dv (\rho \vec u \otimes \vec u) + \varepsilon^\alpha \nabla p + \frac{1}{\varepsilon} \rho \vec u + \rho \nabla \phi & = 0, & \text{in} \ \Omega, \\
        \label{eq:PLEP-3}
        - \Delta \phi & = \rho - M, & \text{in} \ \Omega,
    \end{align}
    where $ p = \rho^\gamma $, $ \gamma > 1 $ is the pressure potential.
\end{subequations} 
Here $ \alpha > 0 $ is a fixed parameter. $ M > 0 $ is the constant background charge. $ 0 < \varepsilon \ll 1 $ is a small parameter. As $ \varepsilon \rightarrow 0 $, system \eqref{sys:PLEP} is said to be in the large friction, almost pressureless regime. 

Furthermore, consider the time scale $ t \simeq 1/\varepsilon $ by introducing the following rescaling of time:
\begin{equation}
    \label{eq:rescale-time}
    t =  \tau/\varepsilon.
\end{equation}
Then the system \eqref{sys:PLEP} can be rewritten as:
\begin{subequations}
    \label{sys:PLEP-rescaled}
    \begin{align}
        \label{eq:rsPLEP-1}
        \dtau \rho + \frac{1}{\varepsilon}\dv (\rho \vec u) & =0, & \text{in} \ \Omega , \\
        \label{eq:rsPLEP-2} 
        \dtau (\rho \vec u) + \frac{1}{\varepsilon}\dv (\rho \vec u \otimes \vec u) + \frac{1}{\varepsilon^{1-\alpha}} \nabla p + \frac{1}{\varepsilon^2} \rho \vec u + \frac{1}{\varepsilon}\rho \nabla \phi & = 0, & \text{in} \ \Omega , \\
        \label{eq:rsPLEP-3}
        - \Delta \phi & = \rho - M, & \text{in} \ \Omega .
    \end{align}
\end{subequations}

Formally, one can check that the leading order dynamics of system \eqref{sys:PLEP-rescaled} as $ \varepsilon \rightarrow 0 $ is given by the following hyperbolic-elliptic Keller-Segel (KS) system of consumption type: 
\begin{subequations}
    \label{sys:PLEP-limit}
    \begin{align}
        \label{eq:rsPLEPlt-1}
        \dtau \sigma - \dv (\sigma \nabla \psi) & =0, & \text{in} \ \Omega \times (0,\infty), \\
        % \label{eq:rsPLEPlt-2} 
        % \dtau (\rho \vec u) + \frac{1}{\varepsilon}\dv (\rho \vec u \otimes \vec u) + \frac{1}{\varepsilon^{1-\alpha}} \nabla p + \frac{1}{\varepsilon^2} \rho \vec u \pm \frac{1}{\varepsilon}\rho \nabla \phi & = 0, & \text{in} \ \Omega \times (0,T/\varepsilon), \\
        \label{eq:rsPLEPlt-3}
        - \Delta \psi & = \sigma - M, & \text{in} \ \Omega \times (0,\infty),
    \end{align}
    or equivalently,
    \begin{equation}{\tag{\ref{sys:PLEP-limit}'}}
        \label{eq:rsPLEPlt-limit}
        \dtau \sigma + \vec v \cdot \nabla \sigma + \sigma (\sigma-M) = 0, \quad \text{with} \quad \vec v = - \nabla (-\Delta)^{-1}(\sigma-M).
    \end{equation}
    where formally, $ \sigma = \lim_{\varepsilon\to 0} \rho $ and $ \psi = \lim_{\varepsilon\to 0} \phi $. 
    \end{subequations}
    Here $ \sigma $ is referred to as the bacteria density, and $ \psi $ is referred to as the chemical potential for the nutrition in the literature. System \eqref{sys:PLEP-limit} is the Keller-Segel system of consumption type. A similar system with opposite sign in front of the divergence in \eqref{eq:rsPLEPlt-1} is referred to as the KS system of production type. 

Motivated by \eqref{eq:rsPLEPlt-limit}, we define the Keller-Segel map for system \eqref{sys:PLEP-rescaled} as follows:
\begin{equation}
    \label{def:KS-map}
    \vec v_\rho: \quad \rho \mapsto  - \nabla \phi = - \nabla (-\Delta)^{-1} (\rho -M),
\end{equation}
satisfying 
\begin{equation}
    \label{def:KS-map-2}
     \vec v_\rho + \nabla \phi = 0.
\end{equation}
Furthermore, let 
\begin{equation}
    \label{def:perturbation-u}
    \vec u := \varepsilon \vec v_\rho + \varepsilon^{\alpha} \vec w.
\end{equation}
Then $ (\rho, \vec w) $ satisfies the following system:
\begin{subequations}
    \label{sys:PLEP-perturbation}
    \begin{align}
        \label{eq:ptb-01}
        \dtau \rho + \frac{1}{\varepsilon^{1-\alpha}} \dv(\rho \vec w) + \dv (\rho \vec v_\rho) & = 0, & \text{in} \ \Omega \times (0,\infty), \\
        \label{eq:ptb-02}
        \varepsilon^\alpha \rho \dtau \vec w + \varepsilon^{\alpha -1} \rho \vec u \cdot \nabla \vec w + \frac{1}{\varepsilon^{1-\alpha}}\nabla p(\rho) + \frac{1}{\varepsilon^{2-\alpha}} \rho \vec w & = -  \varepsilon \rho \dtau \vec v_\rho - \rho \vec u \cdot \nabla \vec v_\rho, & \text{in} \ \Omega \times (0,\infty).
    \end{align}
\end{subequations}

\begin{equation}
    \label{initial:000}
    \norm{\varepsilon^{\alpha/2} \vec w_0}{H^3} + \norm{\rho_0-M}{H^3} < \infty. 
\end{equation}

\begin{equation}
    \label{initial:002}
    \int_{\Omega} (\rho_0 - M) \idx = 0.
\end{equation}

\begin{theorem}[Uniform regularity]
    \label{thm:uniform}
    Let $ 0 < \alpha < 2 $.
    Consider system \eqref{sys:PLEP-perturbation} equipped with initial data $ (\rho_0, \vec w_0) $ satisfying \eqref{initial:000}--\eqref{initial:002}. 
    \begin{enumerate}[label = (\roman*), ref= \ref{thm:uniform} (\roman*)]
        \item \label{thm:uniform-1} There exist $ T \in (0,\infty) $ and small enough $ \varepsilon_1 \in (0,1) $ such that, for all $ \varepsilon \in (0,\varepsilon_1) $, the unique classical solution with initial data $ (\rho_0, \vec w_0) $ to system \eqref{sys:PLEP-perturbation} exists for all $ \tau \in (0,T) $. Furthermore, the solution $ (\rho, \vec w) $ satisfies
        \begin{equation}
            \label{est:thm-01}
            \sup_{0\leq \tau \leq T} \bigl( \norm{\varepsilon^{\alpha/2} \vec w(\tau)}{H^3}^2 + \norm{\rho(\tau)-M}{H^3}^2 \bigr) + \int_0^T \bigl( \norm{\varepsilon^{\alpha/2-1}\vec w(s)}{H^3}^2 + \norm{\rho(s)-M}{H^3}^2 \bigr) \,ds \leq \mathcal C_1 < \infty,  
        \end{equation}
        for some $ \mathcal C_1 \in (0,\infty) $ depending only on the initial data. 

        \item \label{thm:uniform-2} If in addition, $ \norm{\nabla \rho_0}{L^4} $ is small enough, there exist $ \mathfrak c_0 $ and small enough $ \varepsilon_2 \in (0,1) $ such that for all $ \varepsilon \in (0,\varepsilon_2) $, the unique classical solution with initial data $ (\rho_0, \vec w_0) $ to system \eqref{sys:PLEP-perturbation} exists for all $ \tau \in (0,\infty) $. Furthermore, the solution $ (\rho, \vec w) $ satisfies
        \begin{equation}
            \label{est:thm-02}
            \begin{aligned}
                & \sup_{0\leq \tau < \infty}  e^{\mathfrak c_0 \tau}  \bigl( \norm{\varepsilon^{\alpha/2} \vec w(\tau)}{H^3}^2 + \norm{\rho(\tau)-M}{H^3}^2  \bigr)  \\
                & \qquad + \int_0^\infty e^{\mathfrak c_0 \tau} \bigl( \norm{\varepsilon^{\alpha/2-1} \vec w(s)}{H^3}^2 + \norm{\rho(s)-M}{H^3}^2  \bigr) \,ds \leq \mathcal C_2 < \infty,
            \end{aligned}
        \end{equation}
        for some $ \mathcal C_2 \in (0,\infty) $ depending only on the initial data. Moreover, there exist $ \mathfrak c_1 \in (0,\infty) $ and $ \mathfrak c_2 \in (0,\infty)$ such that
        \begin{equation}
            \sup_{0\leq \tau < \infty} \bigl( e^{\mathfrak c_1 \tau} \norm{\rho(\tau)-M}{L^\infty} + e^{\mathfrak c_2 \tau} \norm{\nabla\rho(\tau)}{L^4} \bigr) \leq \mathcal C_3 < \infty, 
        \end{equation}
        for some $ \mathcal C_3 \in (0,\infty) $ depending only on the initial data. 
    \end{enumerate}
    
    
\end{theorem}

\begin{equation}
    \label{ass-prior:rho}
    0 < \frac{1}{2} \underline{\rho} \leq \rho(x,t) \leq 2 \overline \rho  < \infty. 
\end{equation}

\begin{equation}
    \label{ene:708}
    \frac{3}{4} \underline \rho \leq \rho(\vec x, \tau) \leq \frac{3}{2} \overline \rho, 
\end{equation}
which is stronger than \eqref{ass-prior:rho} and thus finishes the proof. 


\section{Large friction limit}
\label{sec:convergence}

Without loss of generality, let $ T $ be the existence time of $ (\rho, \vec w) $, satisfying the following uniform-in-$\varepsilon $ bound:
\begin{equation}
    \label{ene:801}
    \begin{aligned}
        \sup_{0\leq \tau \leq T} \bigl( \norm{\varepsilon^{\alpha/2} \vec w(\tau)}{H^3}^2 + \norm{\rho(\tau)-M}{H^3}^2 \bigr) + \int_0^T \bigl( \norm{\varepsilon^{\alpha/2-1}\vec w(s)}{H^3}^2 + \norm{\rho(s)-M}{H^3}^2 \bigr) \,ds < \infty. 
    \end{aligned}
\end{equation}

\begin{theorem}[Large friction limit]
    \label{thm:limit}
    Under the same assumptions as in theorem \ref{thm:uniform}, there exists a $ \sigma \in L^\infty(0,T;H^2)\cap L^2(0,T;H^3) $ with $ \dtau \sigma \in L^2(0,T;H^2) $, such that as $ \varepsilon \rightarrow 0 $, one has that 
\begin{align}
    \label{thm:limit-1}
    \rho - M  \stackrel{*}{\rightharpoonup} & ~ \sigma - M && \text{in} ~ L^\infty(0,T;H^3);\\
    \label{thm:limit-2}
    \rho - M {\rightharpoonup} & ~ \sigma - M && \text{in} ~ L^2(0,T;H^3); \\
    \label{thm:limit-3}
    \rho - M \rightarrow & ~ \sigma - M && \text{in} ~ C(0,T;H^2_\mrm{loc}); \\
    \label{thm:limit-4}
    \dtau \rho \rightharpoonup & ~ \dtau \sigma && \text{in} ~ L^2(0,T;H^2).
\end{align}
Here $ T \in (0,\infty] $ is the existence time of $ (\rho,\vec w) $ from theorem \ref{thm:uniform}. Moreover, $ \sigma $ solves the limiting system \eqref{sys:PLEP-limit}, and $ \sigma(\tau = 0) = \rho_0 $. 
    
\end{theorem}

\begin{proof}
    The proof of theorem \ref{thm:limit} is the consequence of section \ref{sec:convergence}.
\end{proof}

\end{document}
